\begin{figbox}{34.8}
\begin{tikzpicture}[x=0.78cm,y=0.78cm,line join=round]
  % depth direction
  \def\wx{0.50}\def\wy{0.30}
  % ---- left polarizer ----
  \fill[gray!12] (-5.05,1.03) -- (-4.15,1.57) -- (-4.15,-1.03) -- (-5.05,-1.57) -- cycle;
  \fill[pattern=north east lines,pattern color=black!55]
     (-5.05,1.03) -- (-4.15,1.57) -- (-4.15,-1.03) -- (-5.05,-1.57) -- cycle;
  \draw[thickline] (-5.05,1.03) -- (-4.15,1.57) -- (-4.15,-1.03) -- (-5.05,-1.57) -- cycle;
  % ---- left PI-coated plate ----
  \fill[gray!55] (-4.12,-1.06) rectangle (-4.02,1.06);
  \draw[thickline] (-4.12,-1.06) -- (-4.12,1.06);
  \draw[thickline] (-4.02,-1.06) -- (-4.02,1.06);
  % ---- LC molecules (twist from vertical to into-the-page) ----
  \foreach \x/\a in {-3.6/0,-2.75/11,-1.9/23,-1.05/34,-0.2/45,0.65/56,1.5/68,2.35/79,3.2/90}{
    \pgfmathsetmacro{\dx}{0.34*sin(\a)}
    \pgfmathsetmacro{\dy}{0.58*cos(\a)+0.18*sin(\a)}
    \draw[thickline] (\x-\dx,-\dy) -- (\x+\dx,\dy);
  }
  % ---- right PI-coated plate ----
  \fill[gray!55] (3.82,-1.06) rectangle (3.92,1.06);
  \draw[thickline] (3.82,-1.06) -- (3.82,1.06);
  \draw[thickline] (3.92,-1.06) -- (3.92,1.06);
  % ---- right analyser ----
  \fill[gray!12] (4.15,1.57) -- (5.05,1.03) -- (5.05,-1.57) -- (4.15,-1.03) -- cycle;
  \fill[pattern=north east lines,pattern color=black!55]
     (4.15,1.57) -- (5.05,1.03) -- (5.05,-1.57) -- (4.15,-1.03) -- cycle;
  \draw[thickline] (4.15,1.57) -- (5.05,1.03) -- (5.05,-1.57) -- (4.15,-1.03) -- cycle;
  % ---- labels ----
  \node[above,font=\small] at (-2.0,1.85) {LC molecules};
  \node[below,font=\small] at (-5.35,-1.70) {polarizer};
  \node[below,font=\small] at (-3.45,-1.62) {PI film};
  \node[below,font=\small] at (3.30,-1.62) {PI film};
  \node[below,font=\small] at (4.95,-1.70) {analyser};
  % ---- propagation direction ----
  \draw[force] (-6.1,-2.75) -- (4.3,-2.75);
  \node[above,font=\small] at (-0.9,-2.55) {direction of light propagation};
\end{tikzpicture}
\end{figbox}

Fig.~34.8 shows schematically the twisted alignment of the liquid-crystal molecules in a
$90\degree$ TN-LC cell: the light passes successively through the polarizer, the two PI
alignment films and the analyser in the direction indicated.

\begin{figbox}{34.9}
\begin{tikzpicture}[x=1.0cm,y=1.0cm,line join=round]
  % box (cell) in perspective
  \draw[thickline] (0,0) rectangle (3.6,2.2);
  \draw[thickline] (0,2.2) -- (0.9,2.75) -- (4.5,2.75) -- (3.6,2.2);
  \draw[thickline] (3.6,0) -- (4.5,0.55) -- (4.5,2.75);
  \draw[dashedline] (0,0) -- (0.9,0.55) -- (4.5,0.55);
  \draw[dashedline] (0.9,0.55) -- (0.9,2.75);
  % electric field inside
  \foreach \y in {0.55,1.10,1.65}{
    \draw[force] (0.65,\y) -- (2.95,\y);
  }
  \node[above,font=\small] at (1.8,1.62) {$E$};
  % z axis
  \draw[axis] (4.95,1.0) -- (4.95,2.6) node[right,font=\small] {$z$};
  % polarizer
  \draw[thickline] (-0.12,0.0) -- (-0.12,2.2);
  \draw[thickline] (-0.32,0.0) -- (-0.32,2.2);
  \node[below,font=\small] at (-0.22,-0.15) {polarizer};
  % analyser
  \draw[thickline] (3.85,0.42) -- (3.85,2.62);
  \draw[thickline] (4.05,0.42) -- (4.05,2.62);
  \node[below,font=\small] at (3.95,-0.05) {analyser};
  % top label
  \node[above,font=\small] at (2.25,2.9) {$90\degree$ TN-LC};
  % NB mode arrow
  \draw[force] (-1.55,1.1) -- (-0.45,1.1);
  \node[left,font=\small] at (-1.6,1.1) {NB mode};
\end{tikzpicture}
\end{figbox}

Fig.~34.9: the NB-mode setting of the $90\degree$ TN cell.

\paragraph{II. Experimental procedure}

Set the ``$90\degree$ TN-LC'' cell to the ``normally black'' mode and place it between two
polarizers whose polarization directions are parallel to each other. Using a function
generator, apply a $100\unitv{Hz}$ square-wave voltage to the ITO electrodes on the two sides
of the liquid-crystal cell. Increase the voltage $V_{\mathrm{rms}}$ gradually from $0$ to
$7.2\unitv{V}$.

\begin{itemize}
\item Near the critical points, take several additional sets of data if necessary.
\end{itemize}

\begin{center}
\fbox{\parbox{0.86\linewidth}{\textbf{Problem B-(1)}\quad
Measure the NB $90\degree$ TN-LC cell, tabulate the data and plot the electro-optic
characteristic curve ($I$ versus $V_{\mathrm{rms}}$); determine the switching slope $\gamma$,
$\left(\gamma=\dfrac{V_{90}-V_{10}}{V_{10}}\right)$.}}
\end{center}

\begin{center}
\fbox{\parbox{0.86\linewidth}{\textbf{Problem B-(2)}\quad
Determine the critical voltage $V_c$ of the NB $90\degree$ TN-LC cell. Show clearly on your
graph the method by which you determine the value of $V_c$. Hint: when the applied voltage
exceeds the critical voltage, the transmittance of the light rises steeply.}}
\end{center}

\subsection*{Part C\quad Optical properties of nematic liquid crystals}

\paragraph{Electro-optic switching characteristics of a parallel-aligned liquid-crystal cell}

\paragraph{I. Introduction}

\noindent\textit{Uniformly parallel-aligned liquid-crystal cell (Fig.~34.10)}

A uniformly parallel-aligned liquid-crystal cell is a cell in which the dominant directions of
the front PI alignment film and of the rear PI alignment film are parallel to each other
(Fig.~34.10). If the linear polarization direction of the light incident through the polarizer
is parallel to the dominant LC direction, then in the liquid crystal the light propagates as
an extraordinary wave, and the light undergoes only a single (pure) phase change.

If the incident polarization direction makes an angle of $45\degree$ with the dominant LC
direction (see Fig.~34.11), then, because the ordinary and the extraordinary waves have
different velocities in the LC, a ``phase retardation'' occurs between these two waves. In
this case, with $\theta=45\degree$ and with the polarization directions of the two polarizers
parallel to each other, the transmittance of the light transmitted through the analyser (that
is, the ratio of the transmitted light intensity to the intensity of the light incident on the
cell) is
\[
T_{\parallel}=\cos^{2}\frac{\delta}{2}.
\]

\begin{figbox}{34.10}
\begin{tikzpicture}[x=0.95cm,y=0.95cm,line join=round]
  % left plate
  \fill[gray!18] (0,0) -- (0.45,0.28) -- (0.45,3.08) -- (0,2.8) -- cycle;
  \draw[thickline] (0,0) -- (0.45,0.28) -- (0.45,3.08) -- (0,2.8) -- cycle;
  % right plate
  \fill[gray!18] (7.0,0.28) -- (7.45,0.0) -- (7.45,2.8) -- (7.0,3.08) -- cycle;
  \draw[thickline] (7.0,0.28) -- (7.45,0.0) -- (7.45,2.8) -- (7.0,3.08) -- cycle;
  % LC molecules: undulating chain
  \draw[thickline,decorate,decoration={coil,aspect=0.55,segment length=6.5mm,amplitude=3.2mm}]
      (0.6,1.55) -- (6.85,1.55);
  % label
  \draw[thin] (4.3,1.85) -- (4.3,2.45);
  \node[above,font=\small] at (4.3,2.45) {LC molecules};
  % captions
  \node[below,font=\small,align=center] at (0.2,-0.15) {PI plate\\[-2pt](ITO $+$ PI)};
\end{tikzpicture}
\end{figbox}

Fig.~34.10: a uniformly parallel-aligned liquid-crystal cell, whose two PI plates
(ITO $+$ PI) are parallel to each other.

The ``phase retardation'' $\delta$ in the above expression can be written as
\[
\delta=2\pi d\,\Delta n(V,\lambda)/\lambda ,
\]
where $d$ is the thickness of the liquid-crystal layer, $\lambda$ is the wavelength of the
light in air, $V$ is the root-mean-square value of the alternating voltage, and
$\Delta n\equiv n_e-n_o$ is the birefringence of the liquid crystal; $\Delta n$ is a function
of $\lambda$ and $V$. Both $\Delta n$ and $\delta$ attain their maximum values when $V$ is
equal to zero; moreover, $\Delta n$ decreases as $V$ increases.

In the general case the transmittance can be expressed as follows:
\begin{align*}
T_{\parallel}&=1-\sin^{2}2\theta\,\sin^{2}\frac{\delta}{2},\\[2pt]
T_{\perp}&=\sin^{2}2\theta\,\sin^{2}\frac{\delta}{2},
\end{align*}
where $\parallel$ and $\perp$ denote respectively the two cases in which the polarization
direction of the analyser is parallel and perpendicular to the polarization direction of the
polarizer.

\paragraph{II. Experimental procedure}

\begin{enumerate}
\item[(1)] Replace the $90\degree$ TN liquid-crystal cell by the uniformly parallel-aligned
liquid-crystal cell.
\item[(2)] With $V=0$, set $\theta=45\degree$ (Fig.~34.11). In the experiment, with the
transmission axes of the polarizer and of the analyser perpendicular to each other, rotate the
liquid-crystal specimen until the transmittance reaches its maximum value ($T_{\perp}$); this
determines $\theta=45\degree$. Record $T_{\perp}$. Then set the transmission axes of the
polarizer and of the analyser parallel to each other and measure $T_{\parallel}$ (again with
the applied voltage $V=0$), so as to obtain the phase retardation $\delta$ and the
birefringence $\Delta n$.
\end{enumerate}

\begin{figbox}{34.11}
\begin{tikzpicture}[x=0.92cm,y=0.92cm,line join=round]
  % beam
  \draw[dashedline] (-0.8,0) -- (7.9,0);
  \draw[force] (7.9,0) -- (8.4,0);
  % polarizer
  \draw[thickline] (0.55,0) circle (0.60);
  \draw[<->,line width=0.5pt] (0.05,0) -- (1.05,0);
  \node[below,font=\small] at (0.55,-0.72) {polarizer};
  % cell
  \begin{scope}[shift={(4.15,0)},rotate=45]
    \fill[gray!15] (-1.45,-0.30) rectangle (1.45,0.30);
    \draw[thickline] (-1.45,-0.30) rectangle (1.45,0.30);
    \draw[line width=1.3pt] (-1.45,0.30) -- (1.45,0.30);
    \draw[line width=1.3pt] (-1.45,-0.30) -- (1.45,-0.30);
    \draw[<->,dashedline] (-1.10,0) -- (1.10,0);
    \node at (1.32,0) {\small $L$};
  \end{scope}
  % angle theta
  \draw[thickline] (4.15,0) ++(0.95,0) arc (0:45:0.95);
  \node[font=\small] at (5.02,0.34) {$\theta$};
  % labels of the cell
  \node[font=\small,align=center] at (2.35,1.85) {uniformly\\[-2pt]parallel-aligned\\[-2pt]LC cell};
  \draw[thin] (2.95,1.45) -- (4.20,0.35);
  \node[font=\small] at (5.95,1.75) {$\delta(V,\lambda)$};
  \draw[thin] (5.55,1.60) -- (4.95,0.95);
  % electrodes
  \node[font=\small] at (1.55,-1.55) {$V_{\mathrm{rms}}$};
  \draw[thin] (1.95,-1.35) -- (2.88,-0.92);
  % analyser
  \draw[thickline] (7.25,0) circle (0.60);
  \draw[<->,line width=0.5pt] (6.75,0) -- (7.75,0);
  \draw[<->,line width=0.5pt] (7.25,-0.5) -- (7.25,0.5);
  \node[above,font=\small] at (7.25,0.65) {for $T_{\parallel}$};
  \node[below,font=\small] at (7.25,-0.65) {for $T_{\perp}$};
\end{tikzpicture}
\end{figbox}

Fig.~34.11: schematic diagram of the experimental setup (the arrow $L$ indicates the dominant
direction).

\begin{center}
\fbox{\parbox{0.88\linewidth}{\textbf{Problem C-(1)}\quad
Given that the wavelength of the laser is $\lambda=650\unitv{nm}$, that the thickness of the
liquid-crystal layer is $d=7.7\unitv{\mu m}$ and that $\Delta n\approx0.25$, use the
experimental data for $T_{\perp}$ and $T_{\parallel}$ obtained above to calculate the precise
values of $\delta$ and $\Delta n$ of this cell at $V=0$.}}
\end{center}

\begin{enumerate}
\item[(3)] As in experiment (1) above, with $\theta=45\degree$ fixed, use a function generator
to apply a $100\unitv{Hz}$ square wave to the ITO electrodes at the two sides of the uniformly
parallel-aligned liquid-crystal specimen. With the transmission directions of the polarizer and
of the analyser set parallel to each other, measure the electro-optic switching curve of
$T_{\parallel}$, that is, $T_{\parallel}$ versus the applied voltage $V_{\mathrm{rms}}$, for
$V_{\mathrm{rms}}$ between $0$ and $7\unitv{V}$.
\end{enumerate}

(Hint: the result of measuring the switching curve of the transmittance $T_{\perp}$ may help to
improve the accuracy of the measurement of $T_{\parallel}$; the measured result for
$T_{\perp}$ is not needed in answering the following questions.)

\begin{itemize}
\item Remember to take several additional sets of data near the critical points when necessary
(especially between $0.5$ and $4.0\unitv{V}$).
\end{itemize}

\begin{center}
\fbox{\parbox{0.88\linewidth}{\textbf{Problem C-(2)}\quad
In the $\theta=45\degree$ configuration, measure the parallel-aligned liquid-crystal cell,
tabulate the data and plot the electro-optic switching curve of $T_{\parallel}$.}}
\end{center}

\begin{center}
\fbox{\parbox{0.88\linewidth}{\textbf{Problem C-(3)}\quad
From the data of the electro-optic switching, find the value of the applied voltage
$V_{\pi}$ at which $\delta=\pi$. Hint: $V_{\pi}$ is the applied voltage for which the phase
retardation is $\delta=\pi$ ($180\degree$). Remember that $\Delta n$ is a function of the
applied voltage and that $\Delta n$ decreases as $V$ increases. In fixing the accurate value
of $V_{\pi}$ you may consider using extrapolation.}}
\end{center}

\subsection*{Solution to Problem 4}

\paragraph{Part A\quad Measurement of the electro-optic characteristics of the laser diode}

\paragraph{Problem A-(1)}

(a) Data: a table of data with the variables and units indicated.

\begin{center}
{\small Table 4.2\quad $J$--$I$ data}\\[3pt]
{\small
\renewcommand{\arraystretch}{1.15}
\begin{tabular}{|c|c|c|c|c|c|c|c|c|c|}
\hline
$I/\mathrm{mA}$ & 9.2 & 15.2 & 19.5 & 21.6 & 22.2 & 22.7 & 23.0 & 23.4 & 23.8\\
\hline
$J/\mathrm{V}$ & 0.00 & 0.01 & 0.02 & 0.03 & 0.05 & 0.06 & 0.09 & 0.12 & 0.30\\
\hline\hline
$I/\mathrm{mA}$ & 24.2 & 24.6 & 25.0 & 25.4 & 25.8 & 26.2 & 26.6 & 27.0 & 27.4\\
\hline
$J/\mathrm{V}$ & 0.66 & 1.02 & 1.41 & 1.88 & 2.23 & 2.64 & 3.04 & 3.36 & 3.78\\
\hline\hline
$I/\mathrm{mA}$ & 27.8 & 28.2 & 28.6 & 29.0 & 29.4 & 29.8 & 30.2 & 30.5 & 31.0\\
\hline
$J/\mathrm{V}$ & 4.12 & 4.48 & 4.79 & 5.13 & 5.44 & 5.72 & 6.05 & 6.25 & 6.55\\
\hline\hline
$I/\mathrm{mA}$ & 31.4 & 31.8 & 32.2 & 32.6 & 33.0 & 33.4 & 33.8 & 34.2 & 34.6\\
\hline
$J/\mathrm{V}$ & 6.75 & 6.99 & 7.22 & 7.40 & 7.60 & 7.78 & 7.93 & 8.07 & 8.14\\
\hline\hline
$I/\mathrm{mA}$ & 35.0 & 35.5 & 36.0 & 36.5 & 37.0 & 37.6 & 38.0 & 38.6 & \\
\hline
$J/\mathrm{V}$ & 8.18 & 8.20 & 8.22 & 8.24 & 8.24 & 8.25 & 8.26 & 8.27 & \\
\hline
\end{tabular}}
\end{center}

Current error: $\pm0.1\unitv{mA}$; voltage error: $\pm0.01\unitv{V}$.

(b) Plotting: choose the abscissa and ordinate scales and units that match the range and the
precision of the experimental data.

(c) Drawing the curve: use the appropriate data and a suitable shape of line.

\begin{itemize}
\item As shown in Fig.~34.12: onset $\rightarrow$ $0$ $\rightarrow$ threshold $\rightarrow$
linear $\rightarrow$ saturation.
\end{itemize}

\paragraph{Problem A-(2)}

(a) Mark the linear region of Fig.~34.12; see Fig.~34.13.

(b) Carry out the error analysis by the method of least squares or visually with a ruler.

(c) $I_m\pm\Delta I_m$: obtain suitable data for $I_m$ and the error ($\pm\Delta I_m$) from the
linear region of the $J$--$I$ curve.

(d) A suitable value of $I_{\mathrm{th}}$ and its error:
\[
I_{\mathrm{th}}=[(21\sim26)\pm(0.01\ \text{or}\ 0.2)]\unitv{mA}.
\]

\begin{figbox}{34.12}
\begin{tikzpicture}[x=0.165cm,y=0.52cm]
  % graph paper
  \draw[gray!45,very thin] (0,0) grid[xstep=1,ystep=0.5] (40,9);
  % axes
  \draw[axis] (0,0) -- (41.5,0);
  \draw[axis] (0,0) -- (0,9.4);
  % ticks
  \foreach \x in {0,5,...,40}{\draw ( \x,0) -- (\x,-0.12); \node[below,font=\scriptsize] at (\x,-0.15) {\x};}
  \foreach \y in {0,1,...,9}{\draw (0,\y) -- (-0.8,\y); \node[left,font=\scriptsize] at (-0.9,\y) {\y};}
  \node[below,font=\small] at (20,-0.75) {current $I/\mathrm{mA}$};
  \node[rotate=90,font=\small] at (-4.2,4.5) {light intensity $J/\mathrm{V}$};
  % data
  \draw[thickline] plot[smooth] coordinates {
    (9.2,0.00) (15.2,0.01) (19.5,0.02) (21.6,0.03) (22.2,0.05) (22.7,0.06)
    (23.0,0.09) (23.4,0.12) (23.8,0.30) (24.2,0.66) (24.6,1.02) (25.0,1.41)
    (25.4,1.88) (25.8,2.23) (26.2,2.64) (26.6,3.04) (27.0,3.36) (27.4,3.78)
    (27.8,4.12) (28.2,4.48) (28.6,4.79) (29.0,5.13) (29.4,5.44) (29.8,5.72)
    (30.2,6.05) (30.5,6.25) (31.0,6.55) (31.4,6.75) (31.8,6.99) (32.2,7.22)
    (32.6,7.40) (33.0,7.60) (33.4,7.78) (33.8,7.93) (34.2,8.07) (34.6,8.14)
    (35.0,8.18) (35.5,8.20) (36.0,8.22) (36.5,8.24) (37.0,8.24) (37.6,8.25)
    (38.0,8.26) (38.6,8.27)};
  \foreach \p in {(9.2,0.00),(21.6,0.03),(24.6,1.02),(27.0,3.36),(29.4,5.44),(31.8,6.99),(34.2,8.07),(38.6,8.27)}
    \fill \p circle (1.1pt);
  % annotation
  \draw[thin] (26.2,5.35) -- (28.35,4.72);
  \node[font=\small] at (25.0,5.5) {$I_m\pm\Delta I_m$};
\end{tikzpicture}
\end{figbox}

Fig.~34.12: the $J$--$I$ relation.

\begin{center}
{\small Table 4.3}\\[3pt]
{\small
\renewcommand{\arraystretch}{1.2}
\begin{tabular}{p{0.40\linewidth}|p{0.40\linewidth}}
\hline
Least-squares fit & Visual estimation with a ruler\\
\hline
error in the graph $0.0x\unitv{mA}$ & error in the graph $0.0x\unitv{mA}$\\
\hline
least-squares method & enlarged scale graph\\
\hline
error analysis & draw 3 curves for the error analysis\\
\hline
\end{tabular}}
\end{center}

\begin{figbox}{34.13}
\begin{tikzpicture}[x=1.02cm,y=0.34cm]
  % graph paper
  \draw[gray!45,very thin] (0,22) grid[xstep=0.25,ystep=0.5] (6,32);
  % axes
  \draw[axis] (0,22) -- (6.4,22);
  \draw[axis] (0,22) -- (0,32.6);
  % ticks
  \foreach \x in {0,0.5,...,6}{\draw (\x,22) -- (\x,21.75); \node[below,font=\scriptsize] at (\x,21.7) {\x};}
  \foreach \y in {22,23,...,32}{\draw (0,\y) -- (-0.10,\y); \node[left,font=\scriptsize] at (-0.12,\y) {\y};}
  \node[below,font=\small] at (3.0,20.9) {light intensity $J/\mathrm{V}$};
  \node[rotate=90,font=\small] at (-0.85,27) {current $I/\mathrm{mA}$};
  % fitted straight line
  \draw[thickline] (0,23.47) -- (6.2,30.35);
  % error bars and data
  \foreach \x/\y in {0.66/24.2,1.02/24.6,1.41/25.0,1.88/25.4,2.23/25.8,2.64/26.2,3.04/26.6,
                     3.36/27.0,3.78/27.4,4.12/27.8,4.48/28.2,4.79/28.6,5.13/29.0,5.44/29.4,
                     5.72/29.8,6.05/30.2}{
    \fill (\x,\y) circle (1.0pt);
    \draw[very thin] (\x,\y-0.12) -- (\x,\y+0.12);
    \draw[very thin] (\x-0.06,\y) -- (\x+0.06,\y);
  }
  % annotations
  \draw[thin] (2.2,31.1) -- (2.62,29.55);
  \node[left,font=\small] at (2.18,31.2) {$y$ error line};
  \draw[thin] (3.55,30.55) -- (3.78,27.4);
  \node[right,font=\small] at (3.6,30.7) {$x$ error line};
\end{tikzpicture}
\end{figbox}

Fig.~34.13: the straight line and the extrapolation.

\paragraph{A1-1}

\begin{itemize}
\item Least-squares method
\end{itemize}

\[
I=mJ+b\quad\longrightarrow\quad b=I_{\mathrm{th}},
\]
for $y=mx+b$.

\begin{center}
{\small Table 4.4}\\[3pt]
{\footnotesize
\renewcommand{\arraystretch}{1.12}
\begin{tabular}{|c|c|c|c|c|c|c|}
\hline
$n$ & $y$: $I/\mathrm{mA}$ & $x$: $J$ & $xy$ & $x^{2}$ & $y(x)=mx+b$ & $[y-y(x)]^{2}$\\
\hline
1 & 23.8 & 0.30 & 7.14 & 0.0900 & 23.7937 & 3.969E-05\\
\hline
2 & 24.2 & 0.66 & 15.972 & 0.4356 & 24.17134 & 0.000821\\
\hline
3 & 24.6 & 1.02 & 25.092 & 1.0404 & 24.54898 & 0.00260\\
\hline
4 & 25.0 & 1.41 & 35.25 & 1.9891 & 24.95809 & 0.00176\\
\hline
5 & 25.4 & 1.88 & 47.752 & 3.5344 & 25.45112 & 0.00261\\
\hline
6 & 25.8 & 2.23 & 57.534 & 4.9729 & 25.81827 & 0.000334\\
\hline
7 & 26.2 & 2.64 & 69.168 & 6.9696 & 26.24836 & 0.00234\\
\hline
8 & 26.6 & 3.04 & 80.864 & 9.2416 & 26.66796 & 0.00462\\
\hline
9 & 27.0 & 3.36 & 90.72 & 11.2896 & 27.00364 & 1.325E-05\\
\hline
10 & 27.4 & 3.78 & 103.572 & 14.2884 & 27.44422 & 0.00196\\
\hline
11 & 27.8 & 4.12 & 114.536 & 16.9744 & 27.80088 & 7.733E-07\\
\hline
12 & 28.2 & 4.48 & 126.336 & 20.0704 & 28.17852 & 0.000461\\
\hline
13 & 28.6 & 4.79 & 136.994 & 22.9441 & 28.50371 & 0.00927\\
\hline
 & $\sum y=$ & $\sum x=$ & $\sum xy=$ & $\sum x^{2}=$ & & $\sum[y-y(x)]^{2}$\\
 & 340.6 & 33.71 & 910.93 & 113.840 & & $=0.0268$\\
\hline
\end{tabular}}
\end{center}

\begin{align*}
\Delta&=N\sum x^{2}-\Big(\sum x\Big)^{2}=13(113.840)-(33.71)^{2}=343.6\\[2pt]
m&=\frac{1}{\Delta}\Big(N\sum xy-\sum x\sum y\Big)
 =\frac{13(910.93)-(33.71)(340.6)}{343.556}=1.049\\[2pt]
b&=\frac{1}{\Delta}\Big(\sum x^{2}\sum y-\sum x\sum xy\Big)
 =\frac{(113.840)(340.6)-(33.71)(910.93)}{343.556}=23.48\\[2pt]
\sigma_{y}&=\frac{1}{N-2}\sqrt{\sum\,[y-y(x)]^{2}}=\frac{1}{13-2}\sqrt{0.0268}=0.015\\[2pt]
\sigma&=\sqrt{(\sigma_{y})^{2}+\Big(\frac{\dd y}{\dd x}\sigma_{x}\Big)^{2}}
 =\sqrt{(0.015)^{2}+(1.049\times0.005)^{2}}=0.016\\[2pt]
\sigma_{m}&=\sqrt{\frac{N\sigma^{2}}{\Delta}}=\sqrt{\frac{13\times0.016^{2}}{343.556}}=0.0031\\[2pt]
\sigma_{b}&=\sqrt{\frac{\sum x^{2}}{\Delta}}\;\sigma=0.016\times\sqrt{\frac{113.840}{343.556}}=0.0092
\end{align*}

\[
I_{\mathrm{th}}=(23.48\pm0.01)\unitv{mA}.
\]

\paragraph{A1-2}

\begin{itemize}
\item Visual estimation:
\end{itemize}

\[
I=mJ+b\quad\longrightarrow\quad b=I_{\mathrm{th}},
\]
for $y=mx+b$. Straight line 1: $y=1.00x+23.66$; straight line 2: $y=1.05x+23.48$;
straight line 3: $y=1.13x+23.31$.
\[
I_{\mathrm{th}}(\text{mean})=23.48\unitv{mA},\qquad
\Delta I_{\mathrm{th}}=0.18\unitv{mA},
\]
\[
I_{\mathrm{th}}=(23.5\pm0.2)\unitv{mA}.
\]

\paragraph{Part B\quad Optical properties of the nematic liquid-crystal cell}

\noindent\textbf{Electro-optic switching characteristics of the $90\degree$ TN-LC cell}

\paragraph{Problem B-(1)}

(a) A table of data with the variables and units indicated.

\begin{center}
{\small Table 4.5}\\[3pt]
{\footnotesize
\renewcommand{\arraystretch}{1.08}
\begin{tabular}{|c|c||c|c||c|c|}
\hline
\shortstack{applied\\[-2pt]voltage$/\mathrm{V}$} & \shortstack{light\\[-2pt]intensity$/\mathrm{V}$} & \shortstack{applied\\[-2pt]voltage$/\mathrm{V}$} & \shortstack{light\\[-2pt]intensity$/\mathrm{V}$} & \shortstack{applied\\[-2pt]voltage$/\mathrm{V}$} & \shortstack{light\\[-2pt]intensity$/\mathrm{V}$}\\
\hline
0.00 & 0.00 & 1.78 & 0.21 & 3.41 & 1.39\\
\hline
0.10 & 0.00 & 1.81 & 0.26 & 3.50 & 1.40\\
\hline
0.20 & 0.00 & 1.85 & 0.33 & 3.60 & 1.39\\
\hline
0.30 & 0.00 & 1.90 & 0.44 & 3.70 & 1.40\\
\hline
0.40 & 0.00 & 1.96 & 0.57 & 3.80 & 1.40\\
\hline
0.50 & 0.00 & 2.03 & 0.70 & 4.03 & 1.40\\
\hline
0.60 & 0.00 & 2.08 & 0.80 & 4.22 & 1.40\\
\hline
0.70 & 0.00 & 2.15 & 0.92 & 4.40 & 1.39\\
\hline
0.80 & 0.00 & 2.21 & 1.02 & 4.61 & 1.39\\
\hline
0.90 & 0.00 & 2.28 & 1.10 & 4.78 & 1.40\\
\hline
1.00 & 0.00 & 2.33 & 1.14 & 5.03 & 1.39\\
\hline
1.10 & 0.02 & 2.39 & 1.19 & 5.20 & 1.39\\
\hline
1.20 & 0.04 & 2.44 & 1.22 & 5.39 & 1.38\\
\hline
1.24 & 0.04 & 2.50 & 1.26 & 5.61 & 1.39\\
\hline
1.30 & 0.04 & 2.55 & 1.27 & 5.81 & 1.38\\
\hline
1.34 & 0.03 & 2.60 & 1.29 & 6.02 & 1.38\\
\hline
1.38 & 0.02 & 2.67 & 1.32 & 6.21 & 1.38\\
\hline
1.45 & 0.01 & 2.72 & 1.33 & 6.40 & 1.38\\
\hline
1.48 & 0.01 & 2.85 & 1.36 & 6.63 & 1.38\\
\hline
1.55 & 0.02 & 2.97 & 1.37 & 6.80 & 1.38\\
\hline
1.59 & 0.03 & 3.11 & 1.38 & 7.02 & 1.38\\
\hline
1.64 & 0.05 & 3.20 & 1.39 & 7.20 & 1.38\\
\hline
1.71 & 0.11 & 3.32 & 1.39 & & \\
\hline
\end{tabular}}
\end{center}

(b) Choose suitable scales and units for the abscissa and the ordinate according to the range
and the precision of the experiment.

(c) Measure and plot the $J$--$V_{\mathrm{rms}}$ curve of the light intensity against the
applied voltage (see Fig.~34.14).

\begin{itemize}
\item In the normally black mode the transmitted light intensity is less than
$0.05\unitv{V}$.
\item Before the applied voltage reaches the critical value there is a small light bounce.
\item Once the applied voltage exceeds the critical value, the transmitted light intensity
increases steeply.
\item Once the applied voltage exceeds $3\unitv{V}$, a plateau region appears in the
transmitted light intensity.
\end{itemize}

(d) The value of $\gamma$ and its error.

\begin{itemize}
\item Find the maximum value of the light intensity for applied voltages between $3.0$ and
$7.2\unitv{V}$.
\item Determine the value corresponding to $90\%$ of the maximum light intensity and obtain,
by interpolation, the applied voltage $V_{90}$.
\item Determine the value corresponding to $10\%$ of the maximum light intensity and obtain,
by interpolation, the applied voltage $V_{10}$.
\item The correct value of $\gamma\pm\Delta\gamma$ is $(0.42\sim0.44)\pm0.02$.
\end{itemize}

\begin{figbox}{34.14}
\begin{tikzpicture}[x=0.86cm,y=2.3cm]
  % axes
  \draw[axis] (0,0) -- (7.6,0);
  \draw[axis] (0,0) -- (0,1.58);
  % ticks
  \foreach \x in {0,1,...,7}{\draw (\x,0) -- (\x,-0.02); \node[below,font=\scriptsize] at (\x,-0.02) {\x};}
  \foreach \y in {0,0.5,1.0,1.5}{\draw (0,\y) -- (-0.10,\y); \node[left,font=\scriptsize] at (-0.12,\y) {\y};}
  \node[below,font=\small] at (3.9,-0.28) {applied voltage$/\mathrm{V}$};
  \node[rotate=90,font=\small] at (-0.95,0.79) {light intensity$/\mathrm{V}$};
  % data
  \foreach \x/\y in {0/0,0.1/0,0.2/0,0.3/0,0.4/0,0.5/0,0.6/0,0.7/0,0.8/0,0.9/0,1.0/0,
      1.1/0.02,1.2/0.04,1.24/0.04,1.3/0.04,1.34/0.03,1.38/0.02,1.45/0.01,1.48/0.01,
      1.55/0.02,1.59/0.03,1.64/0.05,1.71/0.11,1.78/0.21,1.81/0.26,1.85/0.33,1.9/0.44,
      1.96/0.57,2.03/0.70,2.08/0.80,2.15/0.92,2.21/1.02,2.28/1.10,2.33/1.14,2.39/1.19,
      2.44/1.22,2.5/1.26,2.55/1.27,2.6/1.29,2.67/1.32,2.72/1.33,2.85/1.36,2.97/1.37,
      3.11/1.38,3.2/1.39,3.32/1.39,3.41/1.39,3.5/1.40,3.6/1.39,3.7/1.40,3.8/1.40,
      4.03/1.40,4.22/1.40,4.4/1.39,4.61/1.39,4.78/1.40,5.03/1.39,5.2/1.39,5.39/1.38,
      5.61/1.39,5.81/1.38,6.02/1.38,6.21/1.38,6.4/1.38,6.63/1.38,6.8/1.38,7.02/1.38,7.2/1.38}
    {\fill (\x,\y) rectangle ++(0.075,0.018);}
  % annotations
  \draw[force] (1.20,0.36) -- (1.20,0.03);
  \node[left,font=\scriptsize] at (1.18,0.40) {$T=0$};
  \draw[force] (1.45,0.36) -- (1.45,0.03);
  \node[right,font=\scriptsize] at (1.47,0.33) {$V_c$};
  \draw[force] (1.76,0.05) -- (1.72,0.13);
  \node[right,font=\scriptsize] at (1.78,0.04) {$V_{10}$};
  \draw[force] (1.95,1.22) -- (2.20,1.30);
  \node[left,font=\scriptsize] at (1.93,1.24) {$V_{90}$};
  \draw[force] (1.15,0.42) -- (1.60,0.10);
  \node[left,font=\scriptsize] at (1.13,0.43) {light bounce};
  \draw[force] (3.55,1.05) -- (3.55,1.36);
  \node[below,font=\scriptsize] at (3.55,1.05) {plateau};
\end{tikzpicture}
\end{figbox}

Fig.~34.14: the $J$--$V_{\mathrm{rms}}$ curve of the $90\degree$ TN-LC cell.

\begin{figbox}{34.15}
\begin{tikzpicture}[x=9.4cm,y=27cm]
  % axes
  \draw[axis] (1.19,0) -- (1.73,0);
  \draw[axis] (1.19,0) -- (1.19,0.175);
  % ticks
  \foreach \x in {1.2,1.3,1.4,1.5,1.6,1.7}{\draw (\x,0) -- (\x,-0.0015); \node[below,font=\scriptsize] at (\x,-0.001) {\x};}
  \foreach \y in {0.00,0.02,...,0.16}{\draw (1.19,\y) -- (1.1925,\y); \node[left,font=\scriptsize] at (1.186,\y) {\y};}
  \node[below,font=\small] at (1.46,-0.016) {applied voltage$/\mathrm{V}$};
  \node[rotate=90,font=\small] at (1.16,0.085) {light intensity$/\mathrm{V}$};
  % data
  \foreach \x/\y in {1.20/0.078,1.23/0.072,1.26/0.066,1.29/0.060,1.32/0.055,1.35/0.049,
      1.38/0.043,1.41/0.032,1.43/0.024,1.45/0.019,1.47/0.021,1.49/0.030,1.51/0.041,
      1.54/0.052,1.57/0.062,1.59/0.072,1.61/0.085,1.63/0.100,1.65/0.118,1.67/0.138,1.69/0.158}
    {\fill (\x,\y) rectangle ++(0.008,0.0035);}
  % Vc
  \draw[force] (1.45,0.055) -- (1.45,0.023);
  \node[right,font=\scriptsize] at (1.46,0.055) {$V_c$};
\end{tikzpicture}
\end{figbox}

Fig.~34.15: the region around $V_c$ on an enlarged scale.

\paragraph{Problem B-(2)}

(a) The value of $V_c$ and its error, $V_c\pm\Delta V_c$.

\begin{itemize}
\item Enlarge the scale of the graph and take more measurement points in the neighbourhood of
$V_c$ (see Fig.~34.15).
\item Determine the value of $V_c$ at which the transmitted light intensity increases
steeply.
\item The correct data for $V_c\pm\Delta V_c$ are $[(1.20\sim1.50)\pm0.01]\unitv{V}$.
\end{itemize}

The data of this graph do not correspond one-to-one to the data of Table 4.5; this graph
merely shows how $V_c$ is obtained.

\paragraph{Part C\quad Optical properties of the nematic liquid-crystal cell}

\noindent\textbf{Electro-optic switching characteristics of the uniformly parallel-aligned
liquid-crystal cell}

\paragraph{Problem C-(1)}

(a) The values of $\delta$ and $\Delta n$ and their errors.

\begin{itemize}
\item Obtain the value of $T_{\parallel}$ and take the average.
\item Determine the order $m$.
\item The correct values of $\delta$ are $15.7\sim18.2$.
\item The correct values of $\Delta n$ are $0.20\sim0.24$.
\end{itemize}

\begin{align*}
T_{\parallel}&=\frac{0.31+0.31+0.31}{3}\unitv{V}=(0.31\pm0.01)\unitv{V},\\[2pt]
T_{\perp}&=\frac{1.04+1.03+1.04}{3}\unitv{V}=(1.04\pm0.01)\unitv{V},
\end{align*}
\[
\tan\frac{\delta}{2}=\pm\sqrt{\frac{T_{\perp}}{T_{\parallel}}}=1.83,
\quad\text{hence}\quad
\delta=4.14+2m\pi\ \ (\text{or}\ -2.14+2m\pi).
\]
\[
\delta=\frac{2\pi d\,\Delta n}{\lambda}=\frac{2\pi\times7.7\times0.25}{0.65}=18.6 .
\]
Taking $m=2$ (or $3$) gives $\delta=16.7$ ($5.32\pi$).
\[
\text{From }\ \delta=\frac{2\pi d\,\Delta n}{\lambda}\ \text{ we obtain }\
\Delta n=\frac{\delta\lambda}{2\pi d}=0.22 .
\]
Permitted values: $\Delta n=(0.20\sim0.24)$.

\begin{itemize}
\item If $\tan\dfrac{\delta}{2}=1.83$, then $\delta$ may be $4.68\pi$ or $6.68\pi$, which is
inconsistent with the data graph of Problem C-(2).
\end{itemize}

\paragraph{Problem C-(2)}

(a) A table of data with the variables and units indicated.

\begin{center}
{\small Table 4.6}\\[3pt]
{\footnotesize
\renewcommand{\arraystretch}{1.08}
\begin{tabular}{|c|c||c|c||c|c|}
\hline
\shortstack{applied\\[-2pt]voltage$/\mathrm{V}$} & \shortstack{light\\[-2pt]intensity$/\mathrm{V}$} & \shortstack{applied\\[-2pt]voltage$/\mathrm{V}$} & \shortstack{light\\[-2pt]intensity$/\mathrm{V}$} & \shortstack{applied\\[-2pt]voltage$/\mathrm{V}$} & \shortstack{light\\[-2pt]intensity$/\mathrm{V}$}\\
\hline
0.00 & 0.30 & 1.08 & 0.25 & 1.50 & 0.13\\
\hline
0.10 & 0.30 & 1.11 & 0.40 & 1.53 & 0.06\\
\hline
0.20 & 0.29 & 1.14 & 0.67 & 1.59 & 0.03\\
\hline
0.30 & 0.29 & 1.17 & 0.93 & 1.62 & 0.05\\
\hline
0.40 & 0.29 & 1.20 & 1.25 & 1.65 & 0.15\\
\hline
0.50 & 0.28 & 1.26 & 1.31 & 1.68 & 0.24\\
\hline
0.60 & 0.26 & 1.29 & 1.36 & 1.71 & 0.34\\
\hline
0.70 & 0.23 & 1.32 & 1.32 & 1.74 & 0.49\\
\hline
0.80 & 0.19 & 1.35 & 1.09 & 1.77 & 0.63\\
\hline
0.90 & 0.09 & 1.38 & 0.85 & 1.80 & 0.78\\
\hline
0.99 & 0.00 & 1.41 & 0.62 & 1.83 & 0.92\\
\hline
1.02 & 0.06 & 1.44 & 0.46 & 1.86 & 1.05\\
\hline
1.05 & 0.16 & 1.47 & 0.29 & 1.89 & 1.19\\
\hline
\end{tabular}}
\end{center}

\begin{center}
{\small Table 4.6 (continued)}\\[3pt]
{\footnotesize
\renewcommand{\arraystretch}{1.08}
\begin{tabular}{|c|c||c|c||c|c|}
\hline
\shortstack{applied\\[-2pt]voltage$/\mathrm{V}$} & \shortstack{light\\[-2pt]intensity$/\mathrm{V}$} & \shortstack{applied\\[-2pt]voltage$/\mathrm{V}$} & \shortstack{light\\[-2pt]intensity$/\mathrm{V}$} & \shortstack{applied\\[-2pt]voltage$/\mathrm{V}$} & \shortstack{light\\[-2pt]intensity$/\mathrm{V}$}\\
\hline
1.92 & 1.27 & 2.79 & 0.26 & 4.20 & 0.11\\
\hline
1.95 & 1.34 & 2.82 & 0.23 & 4.30 & 0.14\\
\hline
1.98 & 1.40 & 2.85 & 0.21 & 4.40 & 0.16\\
\hline
2.01 & 1.47 & 2.88 & 0.18 & 4.50 & 0.19\\
\hline
2.04 & 1.48 & 2.91 & 0.16 & 4.60 & 0.22\\
\hline
2.07 & 1.48 & 2.94 & 0.14 & 4.70 & 0.25\\
\hline
2.10 & 1.48 & 2.97 & 0.12 & 4.80 & 0.28\\
\hline
2.13 & 1.45 & 3.00 & 0.09 & 4.90 & 0.31\\
\hline
2.16 & 1.42 & 3.06 & 0.08 & 5.01 & 0.34\\
\hline
2.19 & 1.38 & 3.09 & 0.06 & 5.11 & 0.37\\
\hline
2.22 & 1.33 & 3.12 & 0.05 & 5.21 & 0.39\\
\hline
2.25 & 1.27 & 3.18 & 0.04 & 5.29 & 0.42\\
\hline
2.28 & 1.20 & 3.21 & 0.03 & 5.39 & 0.44\\
\hline
2.31 & 1.14 & 3.24 & 0.02 & 5.51 & 0.48\\
\hline
2.34 & 1.07 & 3.27 & 0.02 & 5.57 & 0.49\\
\hline
2.37 & 1.00 & 3.30 & 0.01 & 5.70 & 0.52\\
\hline
2.40 & 0.94 & 3.33 & 0.00 & 5.80 & 0.55\\
\hline
2.43 & 0.87 & 3.36 & 0.00 & 5.90 & 0.57\\
\hline
2.46 & 0.79 & 3.39 & 0.00 & 6.01 & 0.60\\
\hline
2.49 & 0.72 & 3.42 & 0.00 & 6.10 & 0.62\\
\hline
2.52 & 0.66 & 3.45 & 0.00 & 6.19 & 0.64\\
\hline
2.55 & 0.61 & 3.48 & 0.00 & 6.30 & 0.66\\
\hline
2.58 & 0.56 & 3.51 & 0.00 & 6.40 & 0.69\\
\hline
2.61 & 0.51 & 3.60 & 0.01 & 6.60 & 0.73\\
\hline
2.64 & 0.46 & 3.70 & 0.02 & 6.70 & 0.74\\
\hline
2.67 & 0.42 & 3.80 & 0.03 & 6.80 & 0.76\\
\hline
2.70 & 0.37 & 3.90 & 0.04 & 7.00 & 0.80\\
\hline
2.73 & 0.33 & 4.00 & 0.07 & 7.20 & 0.83\\
\hline
2.76 & 0.30 & 4.10 & 0.09 & & \\
\hline
\end{tabular}}
\end{center}

(b) Choose suitable scales and units for the abscissa and the ordinate according to the range
and the precision of the experiment.

(c) Measure $T_{\parallel}$ as the applied voltage varies and plot the
$T_{\parallel}$--$V_{\mathrm{rms}}$ curve (see Fig.~34.16).

\begin{itemize}
\item There are 3 minima and 2 sharp maxima of the light.
\item The difference between the two maxima is less than $15\%$.
\item The values at the minima are less than $0.1\unitv{V}$.
\end{itemize}

\begin{figbox}{34.16}
\begin{tikzpicture}[x=0.76cm,y=2.35cm]
  \draw[axis] (0,0) -- (7.5,0);
  \draw[axis] (0,0) -- (0,1.58);
  \foreach \x in {0,1,...,7}{\draw (\x,0) -- (\x,-0.02); \node[below,font=\scriptsize] at (\x,-0.02) {\x};}
  \foreach \y in {0.0,0.5,1.0,1.5}{\draw (0,\y) -- (-0.10,\y); \node[left,font=\scriptsize] at (-0.12,\y) {\y};}
  \node[below,font=\small] at (3.9,-0.28) {applied voltage$/\mathrm{V}$};
  \node[rotate=90,font=\small] at (-0.95,0.79) {light intensity$/\mathrm{V}$};
  \draw[dashedline] plot coordinates {
    (0.00,0.30) (0.20,0.29) (0.40,0.29) (0.60,0.26) (0.80,0.19) (0.90,0.09)
    (0.99,0.00) (1.02,0.06) (1.05,0.16) (1.08,0.25) (1.11,0.40) (1.14,0.67)
    (1.17,0.93) (1.20,1.25) (1.26,1.31) (1.29,1.36) (1.32,1.32) (1.35,1.09)
    (1.38,0.85) (1.41,0.62) (1.44,0.46) (1.47,0.29) (1.50,0.13) (1.53,0.06)
    (1.59,0.03) (1.65,0.15) (1.71,0.34) (1.77,0.63) (1.83,0.92) (1.89,1.19)
    (1.95,1.34) (2.01,1.47) (2.07,1.48) (2.13,1.45) (2.19,1.38) (2.25,1.27)
    (2.31,1.14) (2.37,1.00) (2.46,0.79) (2.55,0.61) (2.64,0.46) (2.73,0.33)
    (2.82,0.23) (2.91,0.16) (3.00,0.09) (3.09,0.06) (3.18,0.04) (3.27,0.02)
    (3.36,0.00) (3.45,0.00) (3.60,0.01) (3.70,0.02) (3.90,0.04) (4.10,0.09)
    (4.30,0.14) (4.50,0.19) (4.70,0.25) (4.90,0.31) (5.11,0.37) (5.29,0.42)
    (5.51,0.48) (5.70,0.52) (5.90,0.57) (6.10,0.62) (6.30,0.66) (6.60,0.73)
    (6.80,0.76) (7.00,0.80) (7.20,0.83)};
  \draw[force] (3.44,0.28) -- (3.44,0.02);
  \node[left,font=\scriptsize] at (3.42,0.30) {$V_{\pi}$};
\end{tikzpicture}
\end{figbox}

Fig.~34.16: the $T_{\parallel}$--$V_{\mathrm{rms}}$ curve.

\paragraph{Problem C-(3)}

The value of $V_{\pi}$ and its error.

\begin{itemize}
\item Enlarge the scale of the graph and take more data in the neighbourhood of $V_{\pi}$
(see Fig.~34.17).
\item Mark the correct value of $V_{\pi}$.
\item Obtain the value of $V_{\pi}$ by interpolation or by rounding off.
\item The correct value of $V_{\pi}$ is $[(3.2\sim3.5)\pm0.01]\unitv{V}$.
\end{itemize}

\begin{figbox}{34.17}
\begin{tikzpicture}[x=4.4cm,y=27cm]
  \draw[axis] (2.98,0) -- (4.03,0);
  \draw[axis] (2.98,0) -- (2.98,0.128);
  \foreach \x in {3.0,3.2,3.4,3.6,3.8,4.0}{\draw (\x,0) -- (\x,-0.0012); \node[below,font=\scriptsize] at (\x,-0.0008) {\x};}
  \foreach \y in {0.00,0.02,0.04,0.06,0.08,0.10,0.12}{\draw (2.98,\y) -- (2.985,\y); \node[left,font=\scriptsize] at (2.975,\y) {\y};}
  \node[below,font=\small] at (3.5,-0.0155) {applied voltage$/\mathrm{V}$};
  \node[rotate=90,font=\small] at (2.94,0.064) {light intensity$/\mathrm{V}$};
  \draw[dashedline] plot coordinates {
    (3.00,0.090) (3.06,0.080) (3.09,0.060) (3.12,0.050) (3.18,0.040) (3.21,0.030)
    (3.24,0.020) (3.27,0.020) (3.30,0.010) (3.33,0.003) (3.36,0.002) (3.39,0.001)
    (3.42,0.001) (3.45,0.001) (3.48,0.002) (3.51,0.003) (3.60,0.010) (3.70,0.020)
    (3.80,0.030) (3.90,0.040) (4.00,0.070)};
  \foreach \x/\y in {3.00/0.090,3.09/0.060,3.18/0.040,3.27/0.020,3.36/0.002,
      3.45/0.001,3.60/0.010,3.70/0.020,3.80/0.030,3.90/0.040}{
    \fill (\x,\y) rectangle ++(0.012,0.0018);}
  \draw[force] (3.45,0.030) -- (3.45,0.004);
  \node[right,font=\scriptsize] at (3.46,0.030) {$V_{\pi}$};
\end{tikzpicture}
\end{figbox}

Fig.~34.17: the region around $V_{\pi}$ on an enlarged scale.
